%%%%

%%%   Самарский государственный технический университет

%%%   Кафедра прикладной математики и информатики

%%%

%%%   Преамбула для статьи в журнал "Вестник Самарского государственного

%%%   технического университета. Серия: «Физико-математические науки»"

%%%   Ориентирована под MikTEx 2.4 for Win32

%%%

%%%   Хотя сам журнал собирается под GNU/Linux на дистрибутиве TeXLive



\documentclass[11pt,a4paper,twoside]{article}



\usepackage[cp1251]{inputenc}

\usepackage[T2A]{fontenc}

\usepackage[english,russian]{babel}

\usepackage{samgtubib}

\usepackage[dvips]{graphicx}

\usepackage[multiple]{footmisc}



\usepackage{samgtu}

\renewcommand{\VestnikYear}{2009} % Год издания Вестника

\renewcommand{\VestnikNumber}{2\,(19)} % Номер Вестника

\renewcommand{\VestnikMonth}{Ноябрь} % Месяц выхода Вестника

\renewcommand{\VestnikData}{01.11.09} % Дата подписания Вестника в печать

\renewcommand{\VestnikUPL}{26,64} % Усл. печ. л.

\renewcommand{\VestnikUIL}{25,73} % Уч.-изд. л.

\renewcommand{\VestnikRegNumber}{479/09} % Регистрационный номер

\renewcommand{\VestnikOrder}{994} % Номер заказа






%\begin{document}

\renewcommand{\firstpage}{132}
\renewcommand{\lastpage}{143}

\udk{517.962.2}

\msc{37C75; 93D05}

\title{Об устойчивости одного класса существенно нелинейных разностных систем}
{Об устойчивости одного класса существенно нелинейных разностных систем}

\author{А.\,А.~Султанбеков}
{С\,у\,л\,т\,а\,н\,б\,е\,к\,о\,в~А.\,А.}

\institute{Санкт-Петербургский государственный университет, \\
198504, Россия, Санкт-Петербург, Петергоф, Университетский просп., 35.}

\email{E-mail}{andrey.sultanbekov@gmail.com}

\workabstract{Изучается проблема устойчивости нулевого решения одного класса
существенно нелинейных разностных систем. 
Доказываются теоремы об устойчивости по неоднородному приближению. 
В~качестве системы нелинейного приближения рассматриваются системы треугольного вида.
Условия, при которых возмущения не нарушают устойчивости нулевого решения, формулируются в виде неравенств,
устанавливающих связь между порядками возмущений и порядками
однородности функций, входящих в систему нелинейного приближения.}

\otherinf{\fullauthor{Андрей Аркадьевич Султанбеков} \position{аспирант} \dept{каф. управления медико"=биологическими системами}}

\keywords{разностная система, функция Ляпунова, асимптотическая устойчивость}

\engtitle{On the stability of a class of essentially nonlinear difference systems}

\engauthor{A.\,A.~Sultanbekov}{S\,u\,l\,t\,a\,n\,b\,e\,k\,o\,v~A.\,A.}

\enginstitute{Saint-Petersburg State University,\\
35, Universitetskii prospekt, Petergof, Saint-Petersburg, 198504, Russia.}


\engabstract{The problem of the zero solution stability for a certain class of essentially nonlinear difference systems is studied. Theorems on the stability by the inhomogeneous approximation are proved. Systems of triangular form are considered as systems of nonlinear approximation. Conditions under which perturbations do not destroy stability of the zero solution are formulated in the form of the inequalities establishing relation between orders of perturbations and homogeneity of functions, entering into the system of nonlinear approximation.}

\otherenginfm{\fullauthor{Andrey A. Sultanbekov}
\position{Postgraduate Student} \dept{Dept. of Medical and Biological Systems Control}}

\engkeywords{difference systems, asymptotic stability, Lyapunov's functions}

\date{14/IX/2011}
\revisiondate{30/V/2012}


\maketitle


\Section[n]{Введение} При построении математической модели
реального процесса важной проблемой является устойчивость.
Один из основных способов исследования устойчивости
для нелинейных систем  обыкновенных дифференциальных уравнений предложил
еще А.\,М.~Ляпунов \cite{sultan:bib1}. В этом способе устойчивость нулевого решения
нелинейной системы сводится к устойчивости линейного приближения
исследуемой системы и условий, накладываемых на нелинейную часть.
Данный подход получил развитие в трудах И.\,Г.~Малкина, Н.\,Н.~Красовского, В.\,И.~Зубова [2--4], где были доказаны теоремы об
устойчивости по нелинейному приближению, и в качестве первого
приближения рассматривались системы с однородными правыми частями.
Также В.\,И.~Зубовым были установлены критерии устойчивости по
обобщенно"=однородному приближению \cite{sultan:bib5}.  Многие из полученных подходов были распространены
на дискретные системы. Результаты в данной области
представлены в работах А.~Халанай, Д.~Векслер, М.\,А.~Скалкиной и
многих других авторов (см.~[6--8] и цитируемую там литературу).

В настоящей работе изучается проблема устойчивости нулевого решения одного класса
существенно нелинейных разностных систем. В качестве систем нелинейного приближения рассматриваются системы треугольного вида.
С помощью дискретного аналога прямого метода Ляпунова доказываются теоремы об
устойчивости по нелинейному приближению. Показывается, что для некоторых
классов систем полученные условия можно ослабить.


\smallskip
\Section[n]{Постановка задачи} Пусть задана система разностных
уравнений
\begin{equation}
x(k+1)=x(k)+f(x(k)),\label{system0}
\end{equation}
где $x(k)=(x_1(k)$,  $x_2(k)$ $\ldots$, $x_n(k))^\top$, $k=0, 1, \ldots$, компоненты вектора $ f(z)$
являются непрерывно дифференцируемыми при всех $ z\in \mathbb E^n$
однородными функциями порядка  $\mu>1$. Таким образом, система
\eqref{system0} имеет нулевое решение.

Наряду с системой \eqref{system0} рассмотрим возмущенную систему
\begin{equation}
x(k+1)=x(k)+f(x(k))+r(k,x(k)).\label{perturbedsystem0}
\end{equation}
Будем считать, что  векторная функция  $r(k, z)$ при $k=0, 1, \ldots$
непрерывна в области $\|z\|<H$, причём для любого $k$ на указанном
множестве справедлива оценка  $\|r(k, z)\|\leq L \, \|z\|^\sigma$.
Здесь $H$, $L$, $\sigma$ "--- положительные постоянные, $\|\cdot\|$ "---
евклидова норма вектора.

Пусть нулевое решение однородной системы дифференциальных
уравнений
$$\dot{z}=f(z)$$
асимптотически устойчиво. Известно \cite[с.~62]{sultan:bib8}, что тогда нулевое
решение системы \eqref{system0} тоже асимптотически устойчиво. В
соответствии с критерием устойчивости по нелинейному
приближению  \cite[с.~70]{sultan:bib8} получаем, что при выполнении неравенства
$\sigma>\mu$ нулевое решение системы \eqref{perturbedsystem0}
также асимптотически устойчиво. Следовательно, возмущения не
нарушают асимптотической устойчивости, если их порядок превышает
порядок однородности функции~$f(z)$.

Предположим теперь, что в правых частях уравнений нелинейного
приближения, кроме однородных членов порядка $\mu$, присутствуют
слагаемые меньшего порядка, т.\,е. вместо системы \eqref{system0}
рассмотрим систему
\begin{equation}
x(k+1)=x(k)+f(x(k))+q(x(k))\label{system01}.
\end{equation}
Здесь векторная функция $q(z)$  непрерывна при $\|z\|<H$ и
удовлетворяет неравенству $\|q(z)\|\leq M \|z\|^\lambda$, $M>0$,
$0<\lambda<\mu$.

Снова считаем, что нулевое решение рассматриваемой системы
асимптотически устойчиво. Покажем, что выполнения прежнего
условия $\sigma>\mu$  уже, вообще говоря, недостаточно для того,
чтобы возмущенная система
\begin{equation}
x(k+1)=x(k)+f(x(k))+q(x(k))+r(k, x(k)) \label{perturbedsystem01}
\end{equation}
имела асимптотически устойчивое нулевое решение.

\begin{example}[1]Рассмотрим систему
\begin{equation}
\begin{cases}
x_1(k+1)=x_1(k)-x_1^\mu(k)+x_1^\alpha(k) x_2^\beta(k),\\
x_2(k+1)=x_2(k)-x_2^\mu(k).\label{example11}
\end{cases}
\end{equation}
Здесь $x_1(k)$, $x_2(k)\in \mathbb R$ при $k=0, 1, \ldots$, $\mu$ "--- рациональное число
с нечётными числителем и
знаменателем, $\mu>1$, $\alpha$ и $\beta$ "--- рациональные числа с
нечётными знаменателями, $\alpha\geq 0$, $\beta>0$,
$\alpha+\beta<\mu$.

Несложно убедиться, что функция Ляпунова $ V(z_1, 
z_2)=z_1^{\gamma_1-\mu+1}+z_2^{\gamma_2-\mu+1}$, где $\gamma_1$,
$\gamma_2$ "--- рациональные числа с нечётными знаменателями и
чётными числителями, $\gamma_1>\mu$, $\gamma_2>\mu$ и
$\gamma_1/\gamma_2>(\mu-\alpha)/\beta$, удовлетворяет всем
условиям дискретного аналога теоремы Ляпунова об асимптотической
устойчивости \cite{sultan:bib8}, т.\,е. нулевое решение системы \eqref{example11}
является асимптотически устойчивым.

Рассмотрим теперь возмущенную систему
\begin{gather}
x_1(k+1)=x_1(k)-x_1^\mu(k)+x_1^\alpha(k) x_2^\beta(k),\\
x_2(k+1)=x_2(k)-x_2^\mu(k)+x_1^\sigma(k),\label{example12}
\end{gather}
где $\sigma$ "--- положительное рациональное число с нечётным
знаменателем.

Исследуем поведение решений системы в неотрицательном квадранте. \linebreak
Нетрудно показать, что существует такое $\overline{H}>0$, что на множестве
$G=\{z=(z_1,z_2)^\top :z\geq0, \|z\|<\overline{H}\}$ векторная функция
$g(z)=({z_1-z_1^\mu+z_1^\alpha z_2^\beta,} $ $ 
z_2-z_2^\mu+z_1^\sigma)^\top $ является неубывающей, т. е. для любых
$\overline{z},$ $ \overline{\overline{z}}\in G$, из
$\overline{z}\geq \overline{\overline{z}}$ следует
$g(\overline{z})\geq g(\overline{\overline{z}})$. Известно [9],
что в области $G$  будет выполняться упорядоченность решений по начальным
данным (если $\overline{x}\geq \overline{\overline{x}}$ и решения $x(k,k_0,\overline{x}),$ $ x(k,k_0,\overline{\overline{x}})\in G$ при $k=k_0,\ldots,\overline{k}$, то $x(k,k_0,
\overline{x})\geq x(k,k_0, \overline{\overline{x}})\geq0$ для всех
$k=k_0,\ldots,\overline{k}$).

Рассмотрим множество $M^+=\{z\in G:v(z)\geq0\}$, где
$v(z)=g(z)-z$, т.\,е. $v(z)=(-z_1^\mu+z_1^\alpha z_2^\beta, 
-z_2^\mu+z_1^\sigma)^\top $. Нетрудно заметить, что если
$x(k, k_0, x_0)\in G$ при $k=k_0, \ldots, k_1$, $x_0\in M^+$, то
$x(k, k_0, x_0)\in M^+$ для всех $k=k_0, \ldots, k_1$. Действительно,
так как из $v(x_0)\geq0$ следует $g(x_0)\geq x_0$, в~силу
упорядоченности решений системы имеем $x(k, k_0, x_0)\leq
x(k, k_0, g(x_0))=g(x(k, k_0, x_0))$ для всех $k=k_0, \ldots, k_1$,
т.\,е. решения, выходящие из точек $x_0\in M^+$, остаются в множестве $M^+$ и тем самым являются неубывающими.

Следовательно, если существует последовательность точек $x^{(m)}\in G$ такая, что
$x^{(m)}\to 0$ при $m\to$ $+\infty$, $x^{(m)}\neq 0$, $v(x^{(m)})\geq
0$, то нулевое решение системы не является асимптотически устойчивым.

Пусть
\begin{equation}
\sigma\leq\mu(\mu-\alpha)/\beta. \label{example13}
\end{equation}
Тогда для любого $\delta>0$ найдутся числа $\theta_1$  и
$\theta_2$, $0<\theta_1<\delta$, $0<\theta_2<\delta$  такие, что $
-\theta_1^\mu+\theta_1^\alpha \theta_2^\beta\geq 0, \,
-\theta_2^\mu+\theta_1^\sigma\geq 0.
$

Таким образом, множество $M^+$ является непустым, и можно выделить
последовательность точек $x^{(m)}\in M^+$ такую, что $x^{(m)}\to
0$ при $m\to+\infty$, $x^{(m)}\neq 0$. Получаем, что при
выполнении неравенства \eqref{example13} нулевое решение системы
\eqref{example12} не может быть асимптотически устойчивым.
В частности, если $\mu=3$, $\alpha=\beta=1$, то условие
\eqref{example13} принимает вид
 $\sigma\leq 6$. Значит, при $\sigma=4$  нулевое решение  системы \eqref{example12}
не является  асимптотически устойчивым.
\end{example}

Рассмотренный пример показывает, что без конкретизации вида правой
части системы \eqref{system01} нельзя указать оценку порядка
возмущений, зависящую только от значений $\mu$ и $\lambda$ и
гарантирующую сохранение асимптотической устойчивости нулевого
решения.

В настоящей статье будем предполагать, что  система нелинейного
приближения \eqref{system01} является треугольной. Определим
условия асимптотической устойчивости нулевого решения для
возмущенных систем, а также рассмотрим классы систем, для которых
возможно уточнение этих условий.

\smallskip
\Section[n]{ Устойчивость по нелинейному треугольному приближению}
Рассмотрим систему разностных уравнений
\begin{equation}
\begin{aligned}
x_1(k+1)&=x_1(k)+f_1(x_1(k))+q(x(k)),\\
x_2(k+1)&=x_2(k)+f_2(x_2(k)),
\end{aligned}
\label{system1}
\end{equation}
где $x_1(k)\in \mathbb E^{n_1}$, $x_2(k)\in \mathbb E^{n_2}$,
$n_1+n_2=n$, $x(k)=(x_1(k), x_2(k))^\top $, элементы векторов $f_1(z_1)$ и
$f_2(z_2)$ "--- непрерывно дифференцируемые однородные функции
порядка $\mu>1$, $z_1\in \mathbb E^{n_1}$, $z_2\in \mathbb E^{n_2}$, $z=(z_1, z_2)^\top $, векторная функция $q(z)$ непрерывна в области
$\|z\|<H$ и удовлетворяет неравенству $\|q(z)\|\leq M\|z_1\|^{\alpha}\|z_2\|^{\beta},$
$H>0$, $M>0$, $\alpha\geq0$, $\beta>0$, $\alpha+\beta<\mu$.
Система \eqref{system1} относится к классу так называемых
треугольных систем \cite{sultan:bib10}.

Будем считать, что нулевые решения однородных систем
дифференциальных уравнений
\begin{equation}
\dot{z}_s=f_s(z_s),\quad s=1, 2,\label{subsystem1}
\end{equation}
асимптотически устойчивы.

Так как системы \eqref{subsystem1} являются асимптотически
устойчивыми,  существуют положительно"=определенные непрерывно
дифференцируемые функции Ляпунова $V_1(z_1)$ и $V_2(z_2)$,
однородные, порядка $\gamma_1-\mu+1$ и ${\gamma_2-\mu+1}$
соответственно, где $\gamma_1$ и $\gamma_2$ "--- любые рациональные
числа, большие $\mu$, с чётными числителями и нечётными
знаменателями, причём $f_s^{*}(z_s)({\partial V_s}/{\partial
z_s})(z_s)$ "--- отрицательно"=определенные функции, $s=1, 2$ \cite{sultan:bib4}.

Согласно теореме 7.1 \cite[с.~87]{sultan:bib8} система \eqref{system1} также имеет асимптотически устойчивое нулевое решение.

Рассмотрим теперь возмущенную систему
\begin{equation}
\begin{aligned}
x_1(k+1)&=x_1(k)+f_1(x_1(k))+q(x(k))+r_1(k,x(k)),\\
x_2(k+1)&=x_2(k)+f_2(x_2(k))+r_2(k,x(k)).\label{perturbedsystem1}
\end{aligned}
\end{equation}
Здесь векторные функции $r_1(k, z)$ и $r_2(k, z)$ при $k=0, 1, \ldots$
непрерывны в области $\|z\|<H$ и удовлетворяют условиям
\begin{equation}
\|r_s(k, z)\|\leq L_s\|z\|^{\sigma},\label{condition1}
\end{equation}
$\sigma>1$, $L_s\geq0$, $s=1, 2$.

\smallskip

\begin{theorem}[1]Если выполнено неравенство
\begin{equation}
\sigma>\mu(\mu-\alpha)/\beta, \label{theorem1}
\end{equation}
то нулевое решение системы \eqref{perturbedsystem1}
 асимптотически
устойчиво\rm.
\end{theorem}

\smallskip

\begin{proof} Для системы \eqref{perturbedsystem1} функцию
Ляпунова строим в виде $ V(z)=V_1(z_1)+V_2(z_2)$, где $V_1(z_1)$ и
$V_2(z_2)$ "--- положительно"=определенные непрерывно
дифференцируемые однородные порядка $\gamma_1-\mu+1$ и
${\gamma_2-\mu+1}$ соответственно функции Ляпунова, $\gamma_1$ и
$\gamma_2$ "--- любые рациональные числа большие $\mu$ с чётными
числителями и нечётными знаменателями. Существует $\tilde{H}>0$
такое, что при $k\geq 0$, $\|x(k)\|<\tilde{H}$ получаем
\begin{multline*}
\Delta
V=\Big(V_1\big(x_1(k)+f_1(x_1(k))+q(x(k))\big)-V_1\big(x_1(k)\big)\Big)+\\+
r_1^{*}(k,x(k))\frac{\partial V_1}{\partial
z_1}\big(x_1(k)+f_1(x_1(k))+q(x(k))+\theta_{1k}r_1(k,x(k))\big)+\\
+\Big(V_2\big(x_2(k+1)+f_2(x_2(k))\big)-V_2\big(x_2(k)\big)\Big)+\\
+r_2^{*}(k,x(k))\frac{\partial V_2}{\partial
z_2}\big(x_2(k)+f_2(x_2(k))+\theta_{2k}r_2(k,x(k))\big)\leq \\
\leq-a_1(\|x_1(k)\|^{\gamma_1}+\|x_2(k)\|^{\gamma_2})+
a_2\Big(\|x_1(k)\|^{\gamma_1-\mu+\alpha}\|x_2(k)\|^{\beta}+\\
+(\|x_1(k)\|^{\alpha}\|x_2(k)\|^{\beta})^{\gamma_1-\mu+1}+
\|x_1(k)\|^{\gamma_1-\mu+\sigma}+\\
+\|x_1(k)\|^{\gamma_1-\mu}\|x_2(k)\|^{\sigma}+(\|x_1(k)\|^{\alpha}\|x_2(k)\|^{\beta})^{\gamma_1-\mu}(\|x_1(k)\|^{\sigma}+\|x_2(k)\|^{\sigma})+\\
+
\|x_1(k)\|^{\sigma(\gamma_1-\mu+1)}+\|x_2(k)\|^{\sigma(\gamma_1-\mu+1)}+\\
+\|x_1(k)\|^{\sigma}\|x_2(k)\|^{\gamma_2-\mu}+\|x_2(k)\|^{\gamma_2-\mu+\sigma}+
\|x_1(k)\|^{\sigma(\gamma_2-\mu+1)}+\|x_2(k)\|^{\sigma(\gamma_2-\mu+1)}\Big),
\end{multline*}
где $a_1$, $a_2$ "---  положительные постоянные,
$\theta_{sk}\in(0, 1)$, $s=1, 2$.

Для существования окрестности положения $z=0$, в которой $\Delta
V$ будет отрицательно"=определенной, достаточно \cite{sultan:bib5}, чтобы выполнялись
следующие неравенства:
\begin{gather}
\notag
\frac{\gamma_1-\mu+\alpha}{\gamma_1}+\frac{\beta}{\gamma_2}>1,
\quad
\frac{\alpha(\gamma_1-\mu+1)}{\gamma_1}+\frac{\beta(\gamma_1-\mu+1)}{\gamma_2}>1,
\quad
\frac{\gamma_1-\mu+\sigma}{\gamma_1}>1,\\ \notag
\frac{\gamma_1-\mu}{\gamma_1}+\frac{\sigma}{\gamma_2}{>}1, 
\frac{(\gamma_1-\mu)\alpha+\sigma}{\gamma_1}+\frac{(\gamma_1-\mu)\beta}{\gamma_2}{>}1, 
\frac{(\gamma_1-\mu)\alpha}{\gamma_1}+\frac{(\gamma_1-\mu)\beta+\sigma}{\gamma_2}{>}1,\\
\frac{(\gamma_1-\mu+1)\sigma}{\gamma_1}>1,\quad
\frac{(\gamma_1-\mu+1)\sigma}{\gamma_2}>1, \quad \frac{\sigma}{\gamma_1}+\frac{(\gamma_2-\mu)}{\gamma_2}>1,\label{theorem12}\\
\notag
\frac{\gamma_2-\mu+\sigma}{\gamma_2}>1, \quad
\frac{\sigma(\gamma_2-\mu+1)}{\gamma_1}>1, \quad
\frac{\sigma(\gamma_2-\mu+1)}{\gamma_2}>1.
\end{gather}
Исключая избыточные неравенства, приходим к следующим ограничениям:
\begin{equation}
\gamma_1/\gamma_2>(\mu-\alpha)/\beta,
\quad\gamma_1/\gamma_2<\sigma/\mu.\label{theorem11}
\end{equation}
Данные условия получаются из первого и
девятого неравенств системы \eqref{theorem12} соответственно.
Все остальные неравенства системы \eqref{theorem12} не
вносят дополнительных ограничений на параметры $\gamma_1$,
$\gamma_2$.

Множество допустимых значений параметров $\gamma_1$, $\gamma_2$ не
является пустым, если выполнено условие \eqref{theorem1}.

Выберем $\gamma_1$, $\gamma_2$, удовлетворяющие неравенствам
\eqref{theorem11}. Тогда для достаточно малых значений
$\|x(k)\|$ и при всех $k\geq 0$ будет справедлива оценка
$$
\Delta V\leq-c_1
\left(\|x_1(k)\|^{\gamma_1}+\|x_2(k)\|^{\gamma_2}\right),
$$
где $c_1$ "--- положительная постоянная. Таким образом, функция
$V(z)$ удовлетворяет  требованиям дискретного аналога теоремы
Ляпунова об асимптотической устойчивости~\cite{sultan:bib8}. 
\end{proof}

\smallskip


\begin{remark}[1] Пример 1 показывает, что сформулированное в
теореме~1 достаточное условие асимптотической устойчивости в~определённом смысле является и необходимым. Действительно, если
при рассматриваемых значениях параметров $\mu$, $\alpha$, $\beta$ и
$\sigma$ неравенство \eqref{theorem1} не выполнено, то можно
построить систему вида \eqref{perturbedsystem1}, нулевое решение которой не
будет асимптотически устойчивым.
\end{remark}

\begin{example}[2]Пусть задано твёрдое тело, вращающееся вокруг
неподвижной точки $O$, расположенной в его центре инерции. Свяжем
с телом оси координат $O\xi\eta\zeta$, которые являются главными
центральными осями этого тела. Уравнения вращательного движения
тела под действием управляющего момента $\mathbf M$ имеют вид
$
\mathbf\Theta \,
\dot{\mathbf\omega}+\mathbf\omega\times\mathbf\Theta\mathbf\omega=\mathbf
M,
$
где $\mathbf \omega=(\omega_1, \omega_2, \omega_3)^\top$ "--- вектор
угловой скорости; $\mathbf\Theta=\mathop{\rm diag}\{A_1, A_2, A_3\}$ "--- тензор
инерции тела; $A_1$, $A_2$, $A_3$ "--- главные центральные моменты
инерции~\cite[c.~413--419]{sultan:bib11}.

Будем считать, что первые два главных центральных момента инерции
тела равны друг другу ($A_1=A_2$), а управляющий момент
определяется по формуле
 $\mathbf M=\left(p_{11}\omega_1^3+p_{12}\omega_2^3, p_{21}\omega_1^3+p_{22}\omega_2^3, p_{33}\omega_3^3\right)^\top$,
 где $p_{ii}<0,$ $i=1, 2, 3,$ $p_{12}p_{21}\leq0$.
Тогда соответствующая система разностных уравнений примет
следующий вид:
\begin{gather*}
A_1\,x_1(k+1)=A_1\,x_1(k)+h\Big(p_{11}x_1^3(k)+p_{12}x_2^3(k)-(A_3-A_1)x_2(k)x_3(k)\Big),\\
A_1\,x_2(k+1)=A_1\,x_2(k)+h\Big(p_{21}x_1^3(k)+p_{22}x_2^3(k)-(A_1-A_3)x_1(k)x_3(k)\Big),\\
A_3\,x_3(k+1)=A_3\,x_3(k)+h p_{33}x_3^3(k).
\end{gather*}
Эта система представляет собой частный случай системы
\eqref{system1} (здесь $\mu=3$, $\alpha=\beta=1$). Следовательно,
нулевое решение системы асимптотически устойчиво.

Рассмотрим теперь возмущенную систему
\begin{gather*}
A_1x_1(k{+}1){=}A_1x_1(k){+}h\Big(p_{11}x_1^3(k){+}p_{12}x_2^3(k){-}(A_3{-}A_1)x_2(k)x_3(k)\Big){+}r_1(k,x(k)),\\
A_1x_2(k{+}1)=A_1x_2(k){+}h\Big(p_{21}x_1^3(k){+}p_{22}x_2^3(k){-}(A_1{-}A_3)x_1(k)x_3(k)\Big){+}r_2(k,x(k)),\\
A_3x_3(k{+}1)=A_3x_3(k){+}h p_{33} x_3^3(k){+}r_3(k,x(k)),
\end{gather*}
где $x(k)=(x_1(k), x_2(k), x_3(k))^\top $, функции $r_i(k, z)$ при $k=0, 1, \ldots$
непрерывны в области $\|z\|<H$ ($H={\rm const}>0$) и удовлетворяют
условиям \mbox{$|r_i(k, z)|\leq N_i\|z\|^{\sigma}$}, $\sigma>0$, $N_i>0$,
$i=1, 2, 3$, $z\in\mathbb E^{3}$.

Применяя теорему 1, получаем, что если $\sigma>6$, то нулевое
решение системы $x(k)\equiv0$ возмущенной системы асимптотически устойчиво.
\end{example}

\smallskip

\begin{newthm}{Следствие} Пусть векторные функции $r_1(k, z)$  и
$r_2(k, z)$ при $k=0, 1, \ldots$ непрерывны в области $\|z\|<H$ и
удовлетворяют неравенствам $$ \|r_s(k, z)\|\leq
L_s\|z_1\|^{\delta_{s1}}\,\|z_2\|^{\delta_{s2}}, $$ где $L_s>0$,
$\delta_{s1}\geq 0$, $\delta_{s2}\geq 0$,
$\delta_{s1}+\delta_{s2}=\sigma>0$, $s=1, 2$. Тогда при
$$\sigma>\big(\mu(\mu-\alpha)-\delta_{22}(\mu-\alpha-\beta)\big)/\beta
$$
нулевое решение системы \eqref{perturbedsystem1} асимптотически
устойчиво.
\end{newthm}

\smallskip

Покажем, что способ исследования устойчивости нелинейных разностных систем по неоднородному треугольному приближению можно использовать и в случае, когда изолированные подсистемы
\eqref{subsystem1} имеют разные порядки однородности.

Пусть в системе \eqref{perturbedsystem1} компоненты векторов $f_1(z_1)$ и
$f_2(z_2)$ являются непрерывно дифференцируемыми однородными
функциями порядка $\mu_1$ и $\mu_2$ соответственно, $\mu_1>1$,
$\mu_2>1$; для непрерывной векторной функции $q(z)$ при $\|z\|<H$
справедлива оценка $\|q(z)\|\leq M\|z_1\|^{\alpha}\|z_2\|^{\beta},$ где $H>0$, $M>0$,
$\alpha\geq 0$, $\beta>0$, а векторная функция $r_2(k, z)$ при
$k=0, 1, \ldots$ непрерывна в области $\|z\|<H$ и удовлетворяет
неравенству $ \|r_2(k, z)\|\leq L_2\|z_1\|^{\delta_{21}}\,
\|z_2\|^{\delta_{22}}$, $L_2>0$, $\delta_{21}>0$, $\delta_{22}\geq
0$. Для простоты будем считать, что $r_1(k, z)\equiv0$.

Заметим, что если $\alpha\geq \mu_1$, то нулевое решение системы
\eqref{perturbedsystem1} асимптотически устойчиво. Для доказательства этого
следует выбирать функцию Ляпунова в виде
$
V(z)=V_1(z_1)+V_2(z_2),
$
где $V_1(z_1)$ и $V_2(z_2)$ "--- положительно"=определенные
непрерывно дифференцируемые однородные порядка ${\gamma_1-\mu_1+1}$
и $\gamma_2-\mu_2+1$ соответственно функции Ляпунова, $\gamma_1$,
$\gamma_2$ "--- рациональные числа с нечётными знаменателями и
чётными числителями,
$\gamma_1>\mu_1$, $\gamma_2>\mu_2$, $\gamma_1/\gamma_2<\delta_{21}/(\mu_2-\delta_{22}),$
причём $f_s^{*}(z_s)({\partial V_s}/{\partial z_s})(z_s)$ "--- отрицательно"=определенные функции, $s=1, 2$ \cite{sultan:bib4}.
Получаем, что для функции Ляпунова $V(z)$ выполнены все условия
дискретного аналога теоремы Ляпунова об асимптотической
устойчивости \cite{sultan:bib8}. Поэтому предположим, что $\alpha<\mu_1$.

\smallskip

\begin{theorem}[2] При выполнении неравенства $ \beta
\delta_{21}>(\mu_1-\alpha)(\mu_2-\delta_{22}) $ нулевое решение
системы \eqref{perturbedsystem1} асимптотически устойчиво\rm.
\end{theorem}

\smallskip
{\it Д\,о\,к\,а\,з\,а\,т\,е\,л\,ь\,с\,т\,в\,о\/} теоремы 2 проводится аналогично доказательству
теоремы~1.

\smallskip

Заметим, что теоремы~1 и~2 будут справедливы и в случае, когда в системе
\eqref{perturbedsystem1}  векторная функция $q$ зависит не только от $z$,
но и от целочисленного аргумента ($q=q(k, z)$), при $k=0, 1, \ldots$ непрерывна в
области $\|z\|<H$ и удовлетворяет неравенству $ \|q(k, z)\|\leq
M\|z_1\|^{\alpha}\,\|z_2\|^{\beta}, $ где $M>0$, $\alpha\geq0$,
$\beta>0$, $\alpha+\beta<\mu$.

Покажем далее, что для некоторых типов функций $q(k, z)$ асимптотическая устойчивость
нулевого решения сохраняется при менее жестких ограничениях на
порядок возмущений по сравнению с условием \eqref{theorem1}.


Пусть $q(k, z)=B(k)g(z_2)$, где матрица $B(k)$
ограничена при $k\geq0$, компоненты вектора $g(z_2)$ являются
непрерывно дифференцируемыми однородными функциями порядка
$\beta$, $1<\beta<\mu$. Тогда возмущенная система имеет вид
\begin{equation}
\begin{aligned}
x_1(k+1)&=x_1(k)+f_1(x_1(k))+B(k)g(x_2(k))+r_1(k,x(k)),\\
x_2(k+1)&=x_2(k)+f_2(x_2(k))+r_2(k,x(k)).\label{perturbedsystem2}
\end{aligned}
\end{equation}
Здесь $f_1(z_1)$ и $f_2(z_2)$ "--- непрерывно дифференцируемые
однородные функции порядка $\mu>1$, векторные функции $r_1(k, z)$ и
$r_2(k, z)$ при $k=0, 1, \ldots$ непрерывны в области $\|z\|<H$ и
удовлетворяют неравенствам \eqref{condition1}. Снова считаем, что
нулевые решения систем \eqref{subsystem1} асимптотически
устойчивы.

В данном случае условие \eqref{theorem1}, при выполнении которого
возмущения не нарушают асимптотическую устойчивость нулевого
решения, принимает вид  $\sigma>\mu^2/\beta$.

Положим
$
C(0)=0$, $C(k)=\sum_{m=0}^{k-1} B(m)$ при $k=1, 2, \ldots$.
Пусть матрица $C(k)$ ограничена для любого  $k=0, 1, \ldots$. В
частности, элементы матрицы $B(k)$ могут представлять собой
периодические последовательности с нулевыми средними значениями.
Тогда полученное ограничение на порядок возмущений можно ослабить
с помощью предложенного в \cite{sultan:bib12} способа построения неавтономных
функций Ляпунова для существенно нелинейных систем.

\smallskip

\begin{theorem}[3] При выполнении неравенства $
\sigma>\mu$ $\max \{1;\, (\mu+1)/(2\beta)\} $
нулевое решение системы \eqref{perturbedsystem2}  асимптотически устойчиво\rm .
\end{theorem}

\smallskip

\begin{proof} Известно \cite{sultan:bib13}, что из асимптотической устойчивости нулевых решений
изолированных подсистем \eqref{subsystem1} следует существование дважды непрерывно
дифференцируемых однородных порядка $\gamma_1-\mu+1$ и
$\gamma_2-\mu+1$ соответственно функций Ляпунова $V_1(z_1)$ и $V_2(z_2)$, удовлетворяющих
требованиям дискретного аналога теоремы Ляпунова об асимптотической устойчивости. В
качестве $\gamma_1$ и $\gamma_2$ можно выбирать любые рациональные
числа, большие $\mu+1$, с~чётными числителями и нечётными знаменателями.

Для системы \eqref{perturbedsystem2} функцию Ляпунова выбираем в виде
$$
{\widetilde V}(k,z)=V_1(z_1)+V_2(z_2)- \left(\frac{\partial
V_1}{\partial z_1}\right)^\top  C(k)g(z_2).
$$
В этом случае найдётся $\breve{H}>0$
такое, что при $k\geq 0$, $\|x(k)\|<\breve{H}$ имеем
\begin{multline*}
\Delta V\leq-b_1\left(\|x_1(k)\|^{\gamma_1}+
\|x_2(k)\|^{\gamma_2}\right)+b_2\Big(\|x_1(k)\|^{\gamma_1-\mu+\sigma}+\|x_1(k)\|^{\gamma_1-\mu}\|x_2(k)\|^{\sigma}+\\
+\|x_1(k)\|^{\gamma_1+\mu-1}+\|x_1(k)\|^{\gamma_1-\mu-1}\|x_2(k)\|^{2\beta}+\|x_1(k)\|^{\gamma_1-\mu-1}
(\|x_1(k)\|^{2\sigma}+\\
+\|x_2(k)\|^{2\sigma})+
+\|x_1(k)\|^{2\mu}\|x_2(k)\|^{\beta(\gamma_1-\mu-1)}+\|x_2(k)\|^{\beta(\gamma_1-\mu+1)}+
\|x_2(k)\|^{\beta(\gamma_1-\mu-1)} \times\\
\times (\|x_1(k)\|^{2\sigma}+\|x_2(k)\|^{2\sigma})+
(\|x_1(k)\|^{2\mu}+\|x_2(k)\|^{2\beta})(\|x_1(k)\|^{\sigma}+\|x_2(k)\|^{\sigma})^{\gamma_1-\mu-1}+\\
+(\|x_1(k)\|^{\sigma}+\|x_2(k)\|^{\sigma})^{\gamma_1-\mu+1}\Big)+
b_3\Big(\|x_2(k)\|^{\gamma_2-\mu}(\|x_1(k)\|^{\sigma}+\|x_2(k)\|^{\sigma})+\\
+(\|x_1(k)\|^{\sigma}+\|x_2(k)\|^{\sigma})^{\gamma_2-\mu+1}\Big)+
b_4\Big(
\|x_1(k)\|^{\gamma_1-\mu}\|x_2(k)\|^{\beta-1+\mu}+\\
+\|x_1(k)\|^{\gamma_1-\mu+\sigma}\|x_2(k)\|^{\beta-1}+\|x_1(k)\|^{\gamma_1-\mu}\|x_2(k)\|^{\beta-1+\sigma}+\\
+\|x_1(k)\|^{\gamma_1-\mu+\sigma(\beta-1)}\|x_2(k)\|^{\mu}+
\|x_1(k)\|^{\gamma_1-\mu}\|x_2(k)\|^{\sigma(\beta-1)+\mu}+\|x_1(k)\|^{\gamma_1-\mu+\sigma\beta}+\\
+\|x_1(k)\|^{\gamma_1-\mu}\|x_2(k)\|^{\sigma\beta}\Big)+b_5\Big(\big(\|x_2(k)\|^{\beta}+(\|x_1(k)\|^{\sigma}+\\
+\|x_2(k)\|^{\sigma})^\beta\big)(\|x_1(k)\|^{\gamma_1-1}+
\|x_1(k)\|^{\gamma_1-\mu-1}\|x_2(k)\|^{\beta}+\\
+\|x_1(k)\|^{\gamma_1-\mu-1}(\|x_1(k)\|^{\sigma}+\|x_2(k)\|^{\sigma})+\|x_1(k)\|^{\mu}\|x_2(k)\|^{\beta(\gamma_1-\mu-1)}+\\
+\|x_2(k)\|^{\beta(\gamma_1-\mu)}+\|x_1(k)\|^{\sigma}\|x_2(k)\|^{\beta(\gamma_1-\mu-1)}+
\|x_2(k)\|^{\beta(\gamma_1-\mu-1)+\sigma}+\\
+(\|x_1(k)\|^{\sigma}+\|x_2(k)\|^{\sigma})^{\gamma_1-\mu-1}(\|x_1(k)\|^{\mu}+\|x_2(k)\|^{\beta})+(\|x_1(k)\|^{\sigma}+\|x_2(k)\|^{\sigma})^{\gamma_1-\mu}\Big).
\end{multline*}

Если параметры $\gamma_1$ и $\gamma_2$ удовлетворяют соотношениям
\begin{equation}
\frac{\gamma_1}{\gamma_2}>\frac{\mu+1}{2\beta},\quad\frac{\gamma_1}{\gamma_2}>\frac{\mu}{\sigma},\quad
\frac{\gamma_1}{\gamma_2}<\frac{\sigma}{\mu}\label{theorem31},
\end{equation}
то для достаточно малых значений $\|x(k)\|$ и
при всех $k\geq 0$ будут справедливы оценки
\begin{gather*}
\widetilde{c}_1
\left(\|x_1(k)\|^{\gamma_1-\mu+1}+\|x_2(k)\|^{\gamma_2-\mu+1}\right)
\leq {\widetilde V}(k,x(k))\leq \\ \hspace{6cm} \leq\widetilde{c}_2
\left(\|x_1(k)\|^{\gamma_1-\mu+1}+\|x_2(k)\|^{\gamma_2-\mu+1}\right),\\
\Delta{\widetilde
V}\leq-\widetilde{c}_3\left(\|x_1(k)\|^{\gamma_1}+\|x_2(k)\|^{\gamma_2}\right), \hspace{3cm}
\end{gather*}
где $\widetilde{c}_1$, $\widetilde{c}_2$, $\widetilde{c}_3$ "---
положительные постоянные.

Множество значений параметров $\gamma_1$, $\gamma_2$, удовлетворяющих соотношениям \eqref{theorem31}, не
является пустым, если выполнено условие теоремы.

Таким образом, функция ${\widetilde V}(k,z)$ удовлетворяет требованиям дискретного аналога теоремы
 Ляпунова об асимптотической устойчивости \cite{sultan:bib8}. 
\end{proof}

\smallskip

\begin{remark}[2] Полученные в теоремах 1, 2, 3 и следствии 1
условия на возмущения для дискретных систем совпадают с условиями,
установленными в работе \cite{sultan:bib14} для непрерывных систем.
\end{remark}

\smallskip

Известно \cite{sultan:bib15}, что если рассматриваемая функция представляет собой
сумму, в которой имеются слагаемые, принимающие в достаточно малой \linebreak
окрестности начала координат неположительные значения, то при
исследовании такой функции на отрицательную определенность
указанные слагаемые можно использовать для ослабления ограничений,
накладываемых на другие слагаемые. Покажем теперь, что для
некоторых классов систем применение данного подхода позволяет
уточнить установленные в теореме 1 условия асимптотической
устойчивости по нелинейному треугольному приближению.


Снова рассмотрим систему \eqref{perturbedsystem1}. Будем считать, что элементы
векторов $f_1(z_1)$ и $f_2(z_2)$ "--- непрерывно  дифференцируемые
при $z_1\in \mathbb E^{n_1}$, $z_2\in \mathbb E^{n_2}$ однородные
функции порядка $\mu>1$, вектор"=функция $q(z)$ непрерывна при
$\|z\|<H$, возмущения $r_1(k, z)$ и $r_2(k, z)$ при $k=0, 1, \ldots$
непрерывны в области $\|z\|<H$ и удовлетворяют неравенствам
$\|r_1(k, z)\|\leq L_1 \, \|z_2\|^{\sigma_1}$, $\|r_2(k, z)\|\hm \leq
L_2 \, \|z_1\|^{\sigma_2}$, $L_1$, $L_2$, $\sigma_1$, $\sigma_2$ "---
положительные постоянные.

Пусть нулевые решения изолированных подсистем  \eqref{subsystem1}
асимптотически устойчивы. Известно \cite{sultan:bib13}, что из асимптотической устойчивости нулевых решений
изолированных подсистем \eqref{subsystem1} следует существование дважды непрерывно
дифференцируемых однородных порядка $\gamma_1-\mu+1$ и
$\gamma_2-\mu+1$ соответственно функций Ляпунова $V_1(z_1)$ и $V_2(z_2)$, удовлетворяющих
требованиям дискретного аналога теоремы Ляпунова об асимптотической устойчивости,
$\gamma_1$ и $\gamma_2$ "--- любые рациональные
числа, б\'ольшие $\mu+1$, с чётными числителями и нечётными
знаменателями. Для системы \eqref{perturbedsystem1} функцию Ляпунова  строим в виде $V(z)=V_1(z_1)+V_2(z_2).$

Предположим, что функцию $V_1(z_1)$ можно выбрать так, чтобы при
\mbox{$\|z\|{<}H$} имела место оценка
$$
\left(\frac{\partial V_1}{\partial
z_1}\Big(z_1)\Big)\right)^\top q(z)\leq -\tilde a
\|z_1\|^{\gamma_1-\mu+\alpha} \|z_2\|^{\beta}.
$$
Здесь $\tilde a>0$, $\alpha\geq 0$, $\beta>0$. Тогда найдётся
$\hat{H}>0$ такое, что при $k\geq 0$, $\|x(k)\|<\hat{H}$ справедливо
соотношение
\begin{multline*}
\Delta V\leq-\overline{c}_1\left(\|x_1(k)\|^{\gamma_1}+
\|x_2(k)\|^{\gamma_2}\right) - \overline{c}
\|x_1(k)\|^{\gamma_1-\mu+ \alpha}\|x_2(k)\|^{\beta}+\\
+\overline{c}_2\Big(\|x_1(k)\|^{\gamma_1-\mu}\|x_2(k)\|^{\sigma_1}+
\|x_1(k)\|^{\gamma_1+\mu-1}+\\
+\|x_1(k)\|^{\alpha(\gamma_1-\mu-1)+2\mu}\|x_2(k)\|^{\beta(\gamma_1-\mu-1)}+
\|x_1(k)\|^{2\mu}\|x_2(k)\|^{\sigma_1(\gamma_1-\mu-1)}+\\
+\|x_1(k)\|^{\gamma_1-\mu-1+2\alpha}\|x_2(k)\|^{2\beta}+(\|x_1(k)\|^{\alpha}\|x_2(k)\|^{\beta})^{\gamma_1-\mu+1}+\\
+\|x_1(k)\|^{2\alpha}\|x_2(k)\|^{\sigma_1(\gamma_1-\mu-1)+2\beta}+
\|x_1(k)\|^{\gamma_1-\mu-1}\|x_2(k)\|^{2\sigma_1}+\\
+\|x_2(k)\|^{\sigma_1(\gamma_1-\mu+1)}+
\|x_1(k)\|^{\alpha(\gamma_1-\mu-1)}\|x_2(k)\|^{\beta(\gamma_1-\mu-1)+2\sigma_1}+\\
+\|x_1(k)\|^{\sigma_2}\|x_2(k)\|^{\gamma_2-\mu}
+\|x_1(k)\|^{\sigma_2(\gamma_2-\mu+1)}\Big),
\end{multline*}
где $\overline{c}_1$, $\overline{c}$, $\overline{c}_2$ "---
неотрицательные постоянные величины, $\theta_{sk}\in(0, 1)$, $s=1, 2$.

Выражение  $-\overline{c} \|x_1(k)\|^{\gamma_1-\mu+
\alpha}\|x_2(k)\|^{\beta}$ неположительно. Поэтому при нахождении
условий отрицательной определенности $\Delta V$  его можно не
учитывать. Если отбросить данное слагаемое, получим, что для
существования окрестности положения $z=0$, в которой $\Delta V$ будет отрицательно"=определенной, достаточно, чтобы выполнялись
следующие неравенства:
\begin{equation*}
\frac{\gamma_1}{\gamma_2}>\frac{\mu}{\sigma_1},\quad
\frac{\gamma_1}{\gamma_2}>\frac{\mu+1-2\alpha}{2\beta},
\quad\frac{\gamma_1}{\gamma_2}<\frac{\sigma_2}{\mu},\quad\gamma_2>\frac{\sigma_2(\mu-1)}{\sigma_2}.
\end{equation*}
Данные соотношения берутся из ограничений, накладываемых на положительные слагаемые
в полученной оценке для приращения функции Ляпунова на решениях системы
\eqref{perturbedsystem1}.

Следовательно, если справедливо неравенство
$$\max \Big\{\cfrac{\mu+1-2\alpha}{2\beta};\cfrac{\mu}{\sigma_1}\Big\}<\cfrac{\sigma_2}{\mu},$$ то
нулевое решение системы \eqref{perturbedsystem1} асимптотически устойчиво.

Пусть имеет место соотношение
$$
\frac{\gamma_1-\mu+\alpha}{\gamma_1}+\frac{\beta}{\gamma_2}<1.
$$
Тогда, применяя лемму 1.2 (см. \cite[с.~9]{sultan:bib15}), получаем, что при
$\sigma_1>\mu\beta/(\mu-\alpha)$ с помощью указанного
слагаемого можно ослабить найденные ограничения  на
значения параметра $\sigma_2$. Нетрудно показать, что в данном
случае асимптотическая устойчивость имеет место при
$\sigma_2>\max\{{\alpha\mu}({\sigma_1-\beta})^{-1}; {\mu(1-\alpha)}{\beta}^{-1}\}$.
Таким образом, справедлива следующая теорема.

\smallskip

\begin{theorem}[4] Для асимптотической устойчивости нулевого
решения системы \eqref{perturbedsystem1} \it достаточно, чтобы при
$\alpha<1$ параметр $\sigma_2$ удовлетворял условиям
\begin{equation*}\sigma_2>
\begin{cases}
\mu^2/\sigma_1& \text{при}\quad \sigma_1\leq \mu\beta/(\mu-\alpha),\\
\alpha\mu/(\sigma_1-\beta)&\text{при}\quad\mu\beta/(\mu-\alpha)<\sigma_1<\beta/(1-\alpha),\\
\mu(1-\alpha)/\beta  &\text{при}\quad \sigma_1>
\beta/(1-\alpha),
\end{cases}
\end{equation*}
а при $\alpha\geq1$ "--- условиям
\begin{equation*}\sigma_2>
\begin{cases}
\mu^2/\sigma_1&\text{при}\quad \sigma_1\leq \mu\beta/(\mu-\alpha),\\
\alpha\mu/(\sigma_1-\beta)&\text{при}\quad \sigma_1>\mu\beta/(\mu-\alpha).
\end{cases}
\end{equation*}
\end{theorem}

\begin{remark}[3]  В случае $\alpha<1$ полученное  в теореме 4
условие на возмущение для дискретных систем является более жестким
по сравнению с условием, установленным в работе \cite{sultan:bib14} для
непрерывного случая. Это связано с тем, что при построении оценки
приращения функции Ляпунова на решениях системы в дискретном случае, в отличие от непрерывного,
 появляются дополнительные слагаемые, которые накладывают свои ограничения.
\end{remark}

\begin{remark}[4] Рассмотренные в статье способы уточнения
условий асимптотической устойчивости по нелинейному неоднородному
приближению применимы и при более общих предположениях
относительно правых частей системы (\ref{system1}). Например,
векторные функции $f_1(z_1)$ и $f_2(z_2)$ могут быть однородными
разного порядка или обобщенно"=однородными, а возмущения могут
удовлетворять условиям $$\|r_1(k, z)\|\leq
L_1\|z_1\|^{\delta_{11}}\|z_2\|^{\delta_{12}},\quad  \|
r_2(k, z)\|\leq L_2\|z_1\|^{\delta_{21}}\|z_2\|^{\delta_{22}},$$
где $L_i>0$, $\delta_{ij}\geq 0$, $\delta_{i1}+\delta_{i2}>0$,
$i, j=1, 2$.
\end{remark}


\begin{thebibliography}{99}

\RBibitem{sultan:bib1}
\by Ляпунов~А.\,М.
\book Общая задача об устойчивости движения
\publaddr М.
\publ ОНТИ
\yr 1935
\totalpages 386
\misctransl
\by Lyapunov~A.\,M.
\book The general problem of the stability of motion.
\publaddr Moscow
\publ ONTI
\yr 1935
\totalpages 386

\RBibitem{sultan:bib2}
\by Малкин~И.\,Г.
\book Теория устойчивости движения
\publaddr М.
\publ Гостехиздат
\yr 1952
\totalpages 432
\misctransl
\by Malkin~I.\,G.
\book Theory of stability of motion
\publaddr Moscow
\publ Gostrkhizdat
\yr 1952
\totalpages 432


\RBibitem{sultan:bib3}
\by Красовский~Н.\,Н.
\book Некоторые задачи теории устойчивости движения
\publaddr М.
\publ Физматгиз
\yr 1959
\totalpages 212
\misctransl
\by Krasovskiy~N.\,N.
\book Some Problems of the Theory of Stability of Motion
\publaddr Moscow
\publ Fizmatgiz
\yr 1959
\totalpages 212

\RBibitem{sultan:bib4}
\by Зубов~В.\,И.
\book Устойчивость движения
\publaddr М.
\publ Высш. шк.
\yr 1973
\totalpages 272
\misctransl
\by Zubov~V.I.
\book Stability of Motion
\publaddr Moscow
\publ Vyssh. Shk.
\yr 1973
\totalpages 272

\RBibitem{sultan:bib5}
\by Зубов~В.\,И.
\book Математические методы исследования систем
автоматического регулирования
\publaddr Л.
\publ Судпромгиз
\yr 1959
\totalpages 324
\misctransl
\by Zubov~V.I.
\book Mathematical Methods for the Study of Automatic Control Systems
\publaddr Leningrad
\publ Sudpromgiz 
\yr 1959
\totalpages 324


\RBibitem{sultan:bib6}
\by Халанай~А., Векcлер~Д.
\book Качественная теория импульсных систем
\publaddr М.
\publ Мир
\yr 1971
\totalpages 312
\misctransl
\by Halanay~A., Vekcler~D.
\book The qualitative theory of sampled-data system
\publaddr Moscow
\publ Mir
\yr 1971
\totalpages 312

\RBibitem{sultan:bib7}
\by Скалкина~М.\,А.
\paper О связи между устойчивостью решений дифференциальных и конечно"=разностных уравнений
\jour Прикл. мат. мех.
\yr 1955
\vol 19
\issue 3
\pages 287--294
\misctransl
\by Skalkina~M.\,A.
\paper On a connection between stability of solutions of differential and finite-difference equations
\jour Prikl. Mat. Mekh.
\vol 19
\issue 3
\yr 1955
\pages 287–294

\RBibitem{sultan:bib8}
\by Александров~А.\,Ю., Жабко~А.\,П.
\book Устойчивость движений дискретных динамических систем 
\publaddr СПб.
\publ СПб. ун-т
\yr 2003
\totalpages 112
\misctransl
\by Aleksandrov~A.\,Yu., Zhabko~A.\,P. 
\book Stability of motion of discrete dynamical systems
\publaddr St. Petersburg
\publ  SPb. Un-t
\yr 2003
\totalpages 112



\RBibitem{sultan:bib9}
\by Козлов~Р.\,И.
\book Теория систем сравнения в методе векторных
функций Ляпунова
\publaddr Новосибирск
\publ Наука
\yr 2001
\totalpages 128
\misctransl
\by Kozlov~R.\,I. 
\book Theory of comparison systems in the method of vector Lyapunov functions
\publ Nauka
\publaddr Novosibirsk
\yr 2001
\totalpages 128

\RBibitem{sultan:bib10}
\by Калитин~Б.\,С.
\paper О принципе сведения для асимптотически
треугольных дифференциальных систем
\jour Прикл. мат. мех.
\yr 2007
\vol 71
\issue 3
\pages 389--400
\transl англ. пер.:~
\by Kalitin~B.\,S. 
\paper On a reduction principle for asymptotically triangular differential systems
\jour J. Appl. Math. Mech.
\vol 71 
\yr 2007
\issue 3
\pages 351–360 

\RBibitem{sultan:bib11}
\by Лурье~А.\,И.
\book Аналитическая механика
\publaddr М.
\publ Физматлит
\yr 1961
\totalpages 824
\misctransl
\by Lur'e~A.\,I.
\book Analytical mechanics
\publ Fizmatlit
\publaddr Moscow 
\yr 1961 
\totalpages 824 

\RBibitem{sultan:bib12}
\by Александров~А.\,Ю.
\book Устойчивость движений неавтономных динамических систем
\publaddr СПб.
\publ СПб. ун-т
\yr 2004
\totalpages 184
\misctransl 
\by Alexandrov~A.\,Yu.
\book Stability of motions of nonautonomous dynamical systems
\publaddr St. Petersburg
\publ  SPb. Un-t
\yr 2004
\totalpages 184


\Bibitem{sultan:bib13}
\by Rosier~L.
\paper Homogeneous Lyapunov function for homogeneous continuous vector field
\jour Systems Control Lett. 
\yr 1992
\vol 19
\issue 6
\pages 467--473

\RBibitem{sultan:bib14}
\by Александров~А.\,Ю., Платонов~А.\,В.
\paper Об устойчивости по нелинейному неоднородному приближению
\jour Матем. заметки
\yr 2011
\vol 90
\issue 6
\pages 803--820
\transl \mbox{англ.} пер.:~
\paper Stability analysis based on nonlinear inhomogeneous approximation
\by  Aleksandrov~A.\,Yu., Platonov~A.\,V.
\jour Math. Notes
\yr 2011
\vol 90
\issue 6
\pages 787--800

\RBibitem{sultan:bib15}
\by Александров~А.\,Ю., Платонов~А.\,В.
\book Устойчивость движений
сложных систем
\publaddr СПб.
\publ СПб. ун-т
\yr 2002
\totalpages 79
\misctransl
\by Aleksandrov~A.\,Yu., Platonov~A.\,V.
\book Stability of complex systems motions
\publaddr St. Petersburg
\publ  SPb. Un-t
\yr 2002
\totalpages 79



\end{thebibliography}

\noindent
\hskip6.5cm \thedate \par\noindent
\hskip6.5cm\therevdate



%\chead[\fancyplain{}{\scriptsize \sl S\,u\,l\,t\,a\,n\,b\,e\,k\,o\,v~A.\,A.}]{\fancyplain{}{\scriptsize \sl Stability of a class of essentially nonlinear
%difference systems}}


\makeengtitlem

\newpage

%\end{document}
